\chapter{Introduction}
Digital imaging research increasingly depends on metadata quality, not only on image visibility. In DICOM repositories, researchers and developers often need to answer questions about protocol consistency, temporal plausibility, dose documentation, device provenance, and cohort comparability before placing confidence in downstream models or quantitative analyses. In practice, these questions are often addressed through ad hoc scripts, manual spot checks in viewers, and repeated one-off exports. This workflow style can be fragile, difficult to reproduce, and inefficient when datasets change over time.

This thesis presents the design and implementation of a full DICOM Metadata Browser as a systems-engineering response to that workflow gap. The central idea is to operationalize metadata work as a persistent pipeline with explicit architecture, storage, validation, and exploration interfaces. Instead of repeatedly rebuilding one-off script outputs, the system provides a stable metadata layer that supports iterative inspection, filtering, quality checks, and export.

The scientific contribution is therefore not a new DICOM parser algorithm, but an integration architecture that unifies heterogeneous file discovery, metadata normalization, persistent, queryable storage, and integrated QA analytics into a single, reproducible artifact. In practical terms, the novelty lies in combining these capabilities under local-first constraints while preserving traceability from dashboard-level anomalies back to concrete series-level metadata records.

The project was developed in the context of a bachelor's thesis and therefore had two explicit constraints. Firstly, the solution had to run reliably in local environments without mandatory external infrastructure. Second, the system had to remain maintainable and understandable within a realistic academic development timeline. These constraints influenced all major design decisions, including the use of Python, pydicom, SQLite, and Flask, the preference for modularized pipeline stages, and the emphasis on interpretable quality checks.

Although the evaluated data originates from medical imaging, the thesis contribution is computer-science oriented: architecture design, pipeline orchestration, data modeling, indexing, concurrency control, validation logic, and reproducible tooling. Clinical metadata examples are used as a demanding application domain for validation, not as the primary research method.

The remainder of this thesis is organized as follows: The first chapters formalize the problem context and derive concrete requirements from observed dataset behavior. The architecture and implementation chapters then explain the processing pipeline in depth, including discovery logic, metadata extraction, private-tag handling, storage design, and web-based exploration. The evaluation chapters present measured runtime and storage behavior on local datasets and analyze quality signals detected by the integrate validation layer. The discussion and conclusion chapters summarize contributions, limitations, and technically grounded future work.
